% Beamer slide deck: how I included `a3_python` in an agential workflow
\documentclass{beamer}
\usetheme{Madrid}
\title{Including a3 in an Agential Workflow}
\author{Automated Demo}
\date{\today}

\begin{document}
\begin{frame}\titlepage\end{frame}

\begin{frame}{Overview}
\begin{itemize}
  \item Goal: use local `a3_python` package inside an agent-style script
  \item Approach: add workspace root to `sys.path`, import `Analyzer`, run `analyze_file()`
\end{itemize}
\end{frame}

\begin{frame}{Files Created}
\begin{itemize}
  \item agential_demo/demo.py -- demo driver (prepends workspace to `sys.path`)
  \item agential_demo/target.py -- small target file to analyze
  \item agential_demo/slides.tex -- this deck
\end{itemize}
\end{frame}

\begin{frame}[fragile]{Key Code (demo.py)}
\begin{verbatim}
workspace_root = Path(__file__).resolve().parents[1]
if str(workspace_root) not in sys.path:
    sys.path.insert(0, str(workspace_root))

from a3_python.analyzer import Analyzer
analyzer = Analyzer(verbose=True)
result = analyzer.analyze_file(target)
print(result.summary())
\end{verbatim}
\end{frame}

\begin{frame}{How It Works}
\begin{enumerate}
  \item Prepend workspace root to `sys.path` so local package imports succeed.
  \item Use `Analyzer` API to compile and symbolically analyze `target.py`.
  \item The demo prints a human-readable `AnalysisResult.summary()`.
\end{enumerate}
\end{frame}

\begin{frame}{Commands to Run}
\begin{verbatim}
python3 agential_demo/demo.py
pdflatex agential_demo/slides.tex
\end{verbatim}
\end{frame}

\end{document}
